\documentclass[11pt,a4paper]{article}

% =========================================================
% راهنمای مطالعه‌ی علوم اعصاب محاسباتی
% نویسنده: Ali Taghizadeh
% برای کامپایل فارسی از XeLaTeX استفاده کنید.
% =========================================================

% ---------- Page Layout ----------
\usepackage[margin=2.3cm]{geometry}

% ---------- Mathematics ----------
\usepackage{amsmath, amssymb, mathtools}

% ---------- Tables and Lists ----------
\usepackage{booktabs}
\usepackage{array}
\usepackage{longtable}
\usepackage{enumitem}

% ---------- Colors and Boxes ----------
\usepackage{xcolor}
\usepackage{tcolorbox}
\tcbuselibrary{skins,breakable}

% ---------- Formatting ----------
\usepackage{titlesec}
\usepackage{fancyhdr}

% ---------- Hyperlinks ----------
\usepackage[hidelinks]{hyperref}

% ---------- Persian Support ----------
% نکته: xepersian باید معمولاً بعد از پکیج‌های اصلی بیاید.
% کامپایل فقط با XeLaTeX
\usepackage{xepersian}

% =========================================================
% Fonts
% =========================================================
% پیشنهاد اصلی: نصب فونت Vazirmatn
% https://github.com/rastikerdar/vazirmatn
%
% نکته‌ی مهم:
% برای متن فارسی و ارقام فارسی نباید از Latin Modern Roman استفاده شود،
% چون کاراکترهای فارسی و ارقام فارسی را ندارد.

\newif\ifpersianfontfound
\persianfontfoundfalse

% ---------- Persian Text Font and Persian Digit Font ----------
\IfFontExistsTF{Vazirmatn}{
	\settextfont{Vazirmatn}
	\setdigitfont{Vazirmatn}
	\persianfontfoundtrue
}{
	\IfFontExistsTF{XB Niloofar}{
		\settextfont{XB Niloofar}
		\setdigitfont{XB Niloofar}
		\persianfontfoundtrue
	}{
		\IfFontExistsTF{Amiri}{
			\settextfont{Amiri}
			\setdigitfont{Amiri}
			\persianfontfoundtrue
		}{
			\IfFontExistsTF{Noto Naskh Arabic}{
				\settextfont{Noto Naskh Arabic}
				\setdigitfont{Noto Naskh Arabic}
				\persianfontfoundtrue
			}{
				\IfFontExistsTF{Noto Sans Arabic}{
					\settextfont{Noto Sans Arabic}
					\setdigitfont{Noto Sans Arabic}
					\persianfontfoundtrue
				}{
					\IfFontExistsTF{DejaVu Serif}{
						\settextfont{DejaVu Serif}
						\setdigitfont{DejaVu Serif}
						\persianfontfoundtrue
					}{
						\IfFontExistsTF{FreeSerif}{
							\settextfont{FreeSerif}
							\setdigitfont{FreeSerif}
							\persianfontfoundtrue
						}{
							% آخرین fallback برای جلوگیری از خطای fontspec.
							% ممکن است حروف فارسی را کامل نمایش ندهد.
							\settextfont{TeX Gyre Termes}
						}
					}
				}
			}
		}
	}
}

% ---------- Latin Text Font ----------
\IfFontExistsTF{TeX Gyre Termes}{
	\setlatintextfont{TeX Gyre Termes}
}{
	\IfFontExistsTF{TeX Gyre Pagella}{
		\setlatintextfont{TeX Gyre Pagella}
	}{
		\IfFontExistsTF{DejaVu Serif}{
			\setlatintextfont{DejaVu Serif}
		}{
			\setlatintextfont{Times New Roman}
		}
	}
}

% =========================================================
% Colors
% =========================================================
\definecolor{navy}{RGB}{43,108,176}
\definecolor{burnt}{RGB}{221,107,32}
\definecolor{leafgreen}{RGB}{47,133,90}
\definecolor{lightgray}{RGB}{245,246,248}
\definecolor{softblue}{RGB}{235,244,255}
\definecolor{darkgray}{RGB}{74,85,104}
\definecolor{midgray}{RGB}{113,128,150}

% =========================================================
% Document Metadata
% =========================================================
\newcommand{\GuideTitle}{از غشای غیرفعال تا شبکه‌های اسپایکی متوازن}
\newcommand{\GuideSource}{یادداشت‌های شخصی بر پایه‌ی مخزن \lr{\texttt{Neuro-AI-Foundations}}}
\newcommand{\GuideAuthor}{\lr{Ali Taghizadeh}}

\hypersetup{
	pdftitle={Computational Neuroscience Study Guide},
	pdfauthor={Ali Taghizadeh},
	pdfsubject={Computational Neuroscience and Spiking Neural Networks},
	pdfkeywords={computational neuroscience, LIF, EIF, AdEx, balanced networks, spiking neural networks}
}

% =========================================================
% Section Styling
% =========================================================
\titleformat{\section}
{\normalfont\Large\bfseries\color{navy}}
{\thesection}
{0.6em}
{}

\titleformat{\subsection}
{\normalfont\large\bfseries\color{darkgray}}
{\thesubsection}
{0.6em}
{}

\titlespacing*{\section}{0pt}{1.4em}{0.7em}
\titlespacing*{\subsection}{0pt}{1.1em}{0.45em}

% =========================================================
% Header / Footer
% =========================================================
\pagestyle{fancy}
\fancyhf{}
\renewcommand{\headrulewidth}{0.4pt}
\renewcommand{\footrulewidth}{0pt}

\fancyhead[R]{\small\textit{راهنمای مطالعه‌ی \lr{Neuro-AI-Foundations}}}
\fancyhead[L]{\small\GuideAuthor\quad|\quad\thepage}

% =========================================================
% List Styling
% =========================================================
\setlist[itemize]{topsep=3pt,itemsep=2pt,parsep=0pt,leftmargin=1.4em}
\setlist[enumerate]{topsep=3pt,itemsep=2pt,parsep=0pt,leftmargin=1.6em}

% =========================================================
% Custom Boxes
% =========================================================
\newtcolorbox{eqbox}[1][]{
	enhanced,
	colback=lightgray,
	colframe=navy,
	boxrule=0.6pt,
	arc=2pt,
	left=6pt,
	right=6pt,
	top=4pt,
	bottom=4pt,
	title={\textbf{معادله‌ی کلیدی: #1}},
	fonttitle=\small\bfseries,
	coltitle=white,
	colbacktitle=navy,
	breakable
}

\newtcolorbox{keybox}[1][نکته‌ی کلیدی]{
	enhanced,
	colback=green!5!white,
	colframe=leafgreen,
	boxrule=0.6pt,
	arc=2pt,
	left=6pt,
	right=6pt,
	top=4pt,
	bottom=4pt,
	title={\textbf{#1}},
	fonttitle=\small\bfseries,
	coltitle=white,
	colbacktitle=leafgreen,
	breakable
}

\newtcolorbox{defbox}[1][تعریف]{
	enhanced,
	colback=orange!5!white,
	colframe=burnt,
	boxrule=0.6pt,
	arc=2pt,
	left=6pt,
	right=6pt,
	top=4pt,
	bottom=4pt,
	title={\textbf{#1}},
	fonttitle=\small\bfseries,
	coltitle=white,
	colbacktitle=burnt,
	breakable
}

\newtcolorbox{overviewbox}{
	enhanced,
	colback=softblue,
	colframe=navy,
	boxrule=0.7pt,
	arc=3pt,
	left=8pt,
	right=8pt,
	top=7pt,
	bottom=7pt,
	width=0.94\textwidth
}

% =========================================================
% Math Shortcuts
% =========================================================
\newcommand{\dd}{\mathrm{d}}
\newcommand{\e}{\mathrm{e}}

% =========================================================
% Title
% =========================================================
\title{}
\author{}
\date{}

\begin{document}
	
	% =========================================================
	% Cover Header
	% =========================================================
	\begin{center}
		\vspace*{-1.2cm}
		
		{\Huge\bfseries\color{navy}\GuideTitle\par}
		

		

		
		\vspace{0.35cm}
		
		{\normalsize\itshape\color{midgray}\GuideSource\par}
		
		\vspace{0.45cm}
		
		{\large\bfseries نویسنده: \GuideAuthor\par}
		

		

		
		\vspace{0.45cm}
		
		{\color{navy}\rule{0.82\textwidth}{0.8pt}}
		
		\vspace{0.55cm}
		
		\begin{overviewbox}
	\centering
	\small
	\textbf{Conceptual progression covered in this guide}\\[6pt]
	
	Passive membrane
	$\;\longrightarrow\;$
	Leaky Integrate-and-Fire, LIF
	$\;\longrightarrow\;$
	Exponential Integrate-and-Fire, EIF
	$\;\longrightarrow\;$
	Adaptive Exponential I\&F, AdEx
	$\;\longrightarrow\;$
	Balanced spiking networks
	
	\vspace{5pt}
	
	{\footnotesize
		\color{darkgray}
		no spikes
		\hspace{0.9cm}
		hard threshold
		\hspace{0.8cm}
		soft nonlinear spike onset
		\hspace{0.8cm}
		adaptation dynamics
		\hspace{0.8cm}
		large-scale network regime
	}
\end{overviewbox}
		
		\vspace{0.45cm}
		
		\begin{tcolorbox}[
			enhanced,
			colback=lightgray,
			colframe=darkgray,
			width=0.94\textwidth,
			arc=3pt,
			boxrule=0.5pt,
			left=8pt,
			right=8pt,
			top=6pt,
			bottom=6pt
			]
			\small
			\textbf{هدف این یادداشت‌ها.}
			این سند یک مسیر ساختاریافته برای مطالعه‌ی مدل‌های اصلی در علوم اعصاب محاسباتی ارائه می‌کند؛
			از معادله‌ی غشای غیرفعال آغاز می‌کند و به‌تدریج به مدل‌های تک‌نورونی تطبیقی و شبکه‌های
			اسپایکی بازگشتی متوازن می‌رسد. تمرکز اصلی روی شهود، فرم ریاضی مدل‌ها، و ارتباط میان
			سازوکارهای زیستی و محاسبات شبکه‌ای در زمینه‌ی هوش مصنوعی است.
		\end{tcolorbox}
		
	\end{center}
	
	\vspace{0.2cm}
	
		\newpage
		
	\tableofcontents
	
	\newpage


% =====================================================================
\section{قرارداد یکاها (ابتدا این بخش را بخوانید)}
% =====================================================================
تمام مدل‌های این پروژه از یک دستگاه یکای یکسان استفاده می‌کنند. پیش از هر چیز باید با این قرارداد راحت شوید، چون دقیقاً همین انتخاب است که باعث می‌شود تمام معادله‌های بعدی از نظر بُعدی سازگار باشند، \emph{بدون} آن‌که به ضرایب تبدیل واحد نیاز داشته باشیم.

\begin{center}
	\begin{tabular}{@{}lll@{}}
		\toprule
		\textbf{کمیت} & \textbf{نماد} & \textbf{یکا} \\
		\midrule
		زمان & $t$ & ms \\
		ولتاژ & $v$ & mV \\
		جریان & $I$ & pA \\
		ظرفیت خازنی & $C_m$ & pF \\
		رسانایی & $g_L$ & nS \\
		مقاومت & $R = 1/g_L$ & G$\Omega$ \\
		\bottomrule
	\end{tabular}
\end{center}

\begin{keybox}[چرا این قرارداد جواب می‌دهد؟]
	داریم:
	$1\,\text{nS} \times 1\,\text{mV} = 1\,\text{pA}$
	و نیز
	$1\,\text{pF} \times 1\,\text{mV/ms} = 1\,\text{pA}$.
	بنابراین معادله‌ای مثل
	$C_m \dot v = -g_L(v-v_{rest}) + I$
	وقتی همه‌چیز با همین یکاها نوشته شود، به هیچ ثابت تبدیل واحدی نیاز ندارد. این همان قراردادی است که در \lr{Brian2} و \lr{NEST} نیز برای این خانواده از مدل‌ها به‌کار می‌رود.
\end{keybox}

% =====================================================================
\section{غشای غیرفعال \lr{(RC)}}
% =====================================================================
این پایه‌ای‌ترین مدل است. غشای نورون، در حداقل توصیف، مانند یک مقاومت و یک خازنِ موازی رفتار می‌کند: بار الکتریکی از طریق مقاومت نشت می‌کند و خازن تغییرات ناگهانی را هموار می‌سازد.

\begin{eqbox}[غشای غیرفعال]
	$$\tau_m \frac{\dd v}{\dd t} = -(v - v_{rest}) + R\,I(t) \qquad\Longleftrightarrow\qquad C_m \frac{\dd v}{\dd t} = -g_L(v - v_{rest}) + I(t)$$
	که در آن
	$\tau_m = C_m/g_L$
	\textbf{ثابت زمانی غشاء} است.
\end{eqbox}

\subsection{پاسخ به جریان پله‌ای (فرم بسته)}
برای یک جریان ثابت $I$ که در $t=0$ وصل می‌شود و با شرط اولیه‌ی $v(0)=v_{rest}$ داریم:
$$v(t) = v_{rest} + \frac{I}{g_L}\left(1 - e^{-t/\tau_m}\right)$$

وقتی $t \to \infty$، خواهیم داشت:
$$v \to v_{rest} + I/g_L = v_{rest} + IR$$

یعنی انحراف حالت ماندگار نسبت به جریان، \emph{خطی} است. این دقیقاً معنای «غیرفعال» بودن است: هیچ رساناییِ وابسته به ولتاژی وجود ندارد، بنابراین سیستم یک فیلتر پایین‌گذر خطی است که فرکانس برش آن را $\tau_m$ تعیین می‌کند.

\begin{keybox}
	در این مدل نه آستانه‌ای وجود دارد و نه بازنشانی \lr{(reset)}. مهم نیست $I$ چقدر بزرگ باشد، $v(t)$ فقط به‌صورت مجانبی به سمت $v_{rest}+IR$ می‌رود و هرگز اسپایک تولید نمی‌کند، چون در خود معادله هیچ سازوکاری برای تولید اسپایک وجود ندارد. دقیقاً به همین دلیل این مدل از نظر آموزشی مهم است: آستانه یک جزئیات فرعی نیست؛ چیزی است که اصلاً مفهوم «نرخ شلیک» را معنادار می‌کند.
\end{keybox}

% =====================================================================
\section{مدل نشت‌دار انباشت و شلیک \lr{(LIF)}}
% =====================================================================
در این مدل، معادله‌ی غشای غیرفعال دقیقاً بدون تغییر حفظ می‌شود و فقط ساده‌ترین قانون ممکن برای اسپایک‌زنی به آن افزوده می‌شود.

\begin{eqbox}[قانون اسپایک در \lr{LIF}]
	$$v(t) \geq v_{th} \;\Longrightarrow\; \text{اسپایک ثبت می‌شود}, \quad v \leftarrow v_{reset}, \quad \text{سپس } v \text{ به‌مدت } t_{ref}\text{ ms نگه داشته می‌شود}$$
\end{eqbox}

\subsection{رئوبیس}
\textbf{رئوبیس} کمینه‌ی جریان ثابت لازم است تا نورون اصلاً بتواند به آستانه برسد؛ یعنی حالت ماندگار دقیقاً به $v_{th}$ برسد:
$$I_{rheo} = g_L\,(v_{th} - v_{rest})$$

اگر جریان از $I_{rheo}$ کمتر باشد، نورون دقیقاً مانند غشای غیرفعال رفتار می‌کند و هرگز اسپایک نمی‌زند. اگر از آن بیشتر باشد، برای همیشه به‌طور تناوبی شلیک می‌کند.

\subsection{منحنی تحلیلی \lr{F-I}}
برای $I > I_{rheo}$، و با فرض ساده‌ساز
$v_{reset}=v_{rest}$،
زمان لازم برای شارژ شدن از مقدار بازنشانی تا آستانه برابر است با:
$$\tau_m \ln\!\left(\dfrac{IR}{IR - (v_{th}-v_{rest})}\right)$$

با افزودن دوره‌ی نارسانایی \lr{refractory period}، فاصله‌ی کامل بین دو اسپایک \lr{inter-spike interval} به‌دست می‌آید و معکوس آن همان نرخ شلیک است:

\begin{eqbox}[منحنی \lr{F-I} برای \lr{LIF}]
	$$f(I) = \left[\, t_{ref} + \tau_m \ln\!\left(\frac{IR}{IR - (v_{th}-v_{rest})}\right)\right]^{-1}, \qquad I > I_{rheo}$$
\end{eqbox}

\begin{keybox}
	وقتی $I \to I_{rheo}^+$، جمله‌ی لگاریتمی واگرا می‌شود و در نتیجه نرخ شلیک به‌صورت پیوسته به صفر میل می‌کند، نه به‌صورت پرشی. از طرف دیگر، وقتی $I \to \infty$، جمله‌ی لگاریتمی به صفر می‌رود و بنابراین
	$f(I) \to 1/t_{ref}$.
	یعنی \textbf{دوره‌ی نارسانایی یک سقف سخت برای نرخ شلیک تعیین می‌کند} و منحنی به‌وضوح در جریان‌های بزرگ به حالت اشباع خم می‌شود.
\end{keybox}

% =====================================================================
\section{انتگرال‌گیری عددی: اویلر در برابر \lr{RK4}}
% =====================================================================
تمام مدل‌های اینجا با گسسته‌سازی زمانی معادله‌ی
$\dot x = f(x,t)$
شبیه‌سازی می‌شوند.

\begin{eqbox}[اویلر رو‌به‌جلو]
	$$x_{n+1} = x_n + \Delta t \cdot f(x_n, t_n)$$
	در هر گام فقط یک بار مشتق ارزیابی می‌شود. این روش ارزان است و \textbf{استاندارد اصلی پروژه برای تمام شبیه‌سازی‌های شبکه‌ای} محسوب می‌شود، با
	$\Delta t \le 0.1$~ms.
\end{eqbox}

\begin{eqbox}[روش رانگ--کوتای کلاسیک مرتبه‌ی چهار \lr{(RK4)}]
	$$k_1 = f(x_n),\quad k_2 = f(x_n+\tfrac{\Delta t}{2}k_1),\quad k_3 = f(x_n+\tfrac{\Delta t}{2}k_2),\quad k_4 = f(x_n+\Delta t\,k_3)$$
	$$x_{n+1} = x_n + \frac{\Delta t}{6}(k_1 + 2k_2 + 2k_3 + k_4)$$
	در هر گام چهار بار مشتق ارزیابی می‌شود؛ بنابراین برای یک $\Delta t$ مشخص، دقت آن بسیار بیشتر است، اما هزینه‌اش هم تقریباً $4\times$ بیشتر است. در این پروژه فقط برای راستی‌آزمایی دقیق مدل تک‌نورونی استفاده می‌شود، نه برای شبکه‌ها.
\end{eqbox}

\begin{keybox}[انتخاب روش چه زمانی واقعاً مهم می‌شود؟]
	برای دینامیک \textbf{خطی} زیرآستانه‌ای، مانند غشای غیرفعال و \lr{LIF}، روش اویلر با
	$\Delta t=0.1$~ms
	تقریباً همان نتیجه‌ی \lr{RK4} را می‌دهد و تفاوت دقت ناچیز است. اما برای دینامیک \textbf{غیرخطی} مانند \lr{EIF} و \lr{AdEx}، خطای موضعی اویلر با افزایش $\Delta t$ سریع‌تر رشد می‌کند؛ در گام‌های درشت‌تر
	$(\Delta t \gtrsim 2~\text{ms})$
	اویلر ممکن است اسپایک‌هایی را که یک مرجع ریزدانه یا \lr{RK4} هنوز تشخیص می‌دهد، از دست بدهد. دقیقاً به همین دلیل شبیه‌سازی‌های شبکه‌ای
	$(1000 \text{ نورون} \times \text{ده‌ها هزار گام زمانی})$
	از اویلر ارزان با
	$\Delta t \le 0.1$~ms
	استفاده می‌کنند؛ جایی که این خطا برای همه‌ی مدل‌های پروژه ناچیز است. در مقابل، \lr{RK4} برای بررسی‌های دقیق روی یک نورون واحد باقی می‌ماند.
\end{keybox}

% =====================================================================
\section{مدل نمایی انباشت و شلیک \lr{(EIF)}}
% =====================================================================
آستانه‌ی سخت در \lr{LIF} بیشتر یک سهولت مدل‌سازی است تا یک سازوکار واقعی. در نورون‌های حقیقی، آغاز اسپایک توسط فعال‌سازی سریع و شدیداً وابسته‌به‌ولتاژ کانال‌های سدیمی
$\text{Na}^+$
ایجاد می‌شود؛ چیزی که با یک شتاب‌گیری نمایی بسیار بهتر از یک تابع پله‌ای تقریب زده می‌شود.

\begin{eqbox}[معادله‌ی غشایی \lr{EIF}]
	$$C_m \frac{\dd v}{\dd t} = -g_L(v - E_L) + g_L \Delta_T \exp\!\left(\frac{v - v_T}{\Delta_T}\right) + I(t)$$
	وقتی $v$ به یک برش عددی
	$v_{peak}$
	برسد، یک اسپایک ثبت می‌شود
	(این مقدار جایگزین بخش صعودیِ واگرای واقعی اسپایک است)؛ سپس
	$v \leftarrow v_{reset}$.
\end{eqbox}

\subsection{نقش \texorpdfstring{$\Delta_T$}{Delta\_T} (عامل شیب آستانه)}
\begin{itemize}[itemsep=2pt]
	\item اگر $\Delta_T \to 0$، جمله‌ی نمایی به یک پله در $v_T$ نزدیک می‌شود و در نتیجه رفتاری نزدیک به آستانه‌ی سختِ \lr{LIF} بازسازی می‌شود.
	\item هرچه $\Delta_T$ بزرگ‌تر باشد، آغاز اسپایک نرم‌تر و تدریجی‌تر می‌شود و خمیدگی بخش صعودی قبل از قله واضح‌تر خواهد بود.
\end{itemize}

\begin{keybox}
	چون جمله‌ی نمایی حتی \emph{پایین‌تر} از $v_T$ نیز ناصفر است، پیش از آستانه هم یک تقویت دپولاریزانِ کوچک ایجاد می‌کند. بنابراین اگر
	$v_T = v_{th}$
	را همسان بگیریم، \textbf{رئوبیس مؤثر \lr{EIF} از \lr{LIF} کمتر است}؛ چون بخشی از کاری که در \lr{LIF} فقط با افزایش جریان انجام می‌شود، در اینجا توسط غیرخطیّت نمایی تأمین می‌شود.
\end{keybox}

\subsection{تحریک‌پذیری نوع اوّل}
هر دو مدل \lr{LIF} و \lr{EIF} دارای \textbf{تحریک‌پذیری نوع اوّل} هستند: منحنی \lr{F-I} آن‌ها از صفر و به‌صورت پیوسته در رئوبیس بالا می‌آید و جهش ناپیوسته به یک نرخ ناصفر ندارد. همچنین تأخیر تا اولین اسپایک وقتی
$I \to I_{rheo}^+$
می‌رود، \textbf{واگرا} می‌شود؛ این همان نشانۀ «کندشدگی بحرانی» نزدیک دوشاخه‌گی‌ای است که شلیک تکراری را ایجاد می‌کند
(نگاه کنید به بخش \lr{\S7}).

% =====================================================================
\section{\lr{AdEx}: مدل نمایی تطبیقی انباشت و شلیک}
% =====================================================================
نه \lr{LIF} و نه \lr{EIF} نمی‌توانند یکی از رفتارهای بسیار رایج نورون‌های قشری را بازتولید کنند: این‌که در طول یک پله‌ی جریان ثابت، \emph{نرخ شلیک خود را تغییر دهند}
(\lr{spike-frequency adaptation})
یا به‌صورت خوشه‌ای و فشرده شلیک کنند
(\lr{bursting}).
هر دو رفتار به یک متغیر حالت دوم نیاز دارند: یک جریان کند به نام $w$ که با شلیک‌ها ساخته می‌شود و دوباره روی $v$ بازخورد می‌دهد.

\begin{eqbox}[معادلات جفت‌شده‌ی \lr{AdEx}]
	$$C_m \frac{\dd v}{\dd t} = -g_L(v - E_L) + g_L \Delta_T \exp\!\left(\frac{v - v_T}{\Delta_T}\right) - w + I(t)$$
	$$\tau_w \frac{\dd w}{\dd t} = a(v - E_L) - w + b\,\tau_w \sum_f \delta(t - t^f)$$
	در هر اسپایک
	$(v \to v_{peak})$
	داریم:
	$\quad v \leftarrow v_{reset}, \qquad w \leftarrow w + b$
\end{eqbox}

\begin{defbox}[چگونه جمله‌ی \lr{$\delta$-function} را بخوانیم؟]
	عبارت
	$b\,\tau_w\sum_f \delta(t-t^f)$
	دقیقاً همان کاری را می‌کند که از روی ظاهرش انتظار می‌رود: در هر زمان اسپایک $t^f$، متغیر $w$ یک \textbf{جهش آنی به اندازه‌ی $b$} دریافت می‌کند. بین اسپایک‌ها، این جمله اصلاً حضور مؤثری ندارد و $w$ با ثابت زمانی $\tau_w$ به سمت
	$a(v - E_L)$
	آرام می‌گیرد. این جمله یک تقریب هموارشده یا فیلترشده نیست؛ در پیاده‌سازی باید جهش مستقیماً اعمال شود.
\end{defbox}

\subsection{هر پارامتر چه چیزی را کنترل می‌کند؟}
\begin{center}
	\begin{tabular}{@{}p{1.8cm}p{2.6cm}p{9.1cm}@{}}
		\toprule
		\textbf{پارامتر} & \textbf{یکا} & \textbf{نقش} \\
		\midrule
		$a$ & nS & کوپلینگ \emph{زیرآستانه‌ای}: این‌که $w$ حتی بدون اسپایک تا چه حد $v$ را دنبال کند. مقادیر منفی $a$ (تسهیل‌گری) همان چیزی هستند که امکان \lr{bursting} را فراهم می‌کنند. \\
		$b$ & pA & افزایش \emph{برانگیخته‌شده با اسپایک} در $w$. هرچه $b$ بزرگ‌تر باشد، سازگاری ناشی از هر اسپایک قوی‌تر است. \\
		$\tau_w$ & ms & ثابت زمانی فروکاهش سازگاری. $\tau_w$ کوچک یعنی بازیابی سریع بین اسپایک‌ها (که می‌تواند خوشه‌های اولیه را ممکن کند)؛ $\tau_w$ بزرگ یعنی سازگاری مدت بیشتری باقی می‌ماند و نرخ شلیک را برای زمان طولانی‌تری سرکوب می‌کند. \\
		$v_{reset}$ & mV & ولتاژی که $v$ پس از اسپایک به آن بازمی‌گردد. اگر به $v_T$ نزدیک یا از آن بالاتر باشد، جمله‌ی نمایی پس از بازنشانی هنوز فعال است و می‌تواند پیش از آن‌که $w$ فرصت جبران پیدا کند، شلیک مجدد سریع ایجاد کند. \\
		\bottomrule
	\end{tabular}
\end{center}

\subsection{چهار رژیم، یک مدل}
\begin{center}
	\begin{tabular}{@{}lllll@{}}
		\toprule
		\textbf{رژیم} & $a$ (nS) & $b$ (pA) & $\tau_w$ (ms) & $v_{reset}$ (mV) \\
		\midrule
		شلیک تونیک & 0 & 0 & 100 & $-58$ \\
		سازگارشونده & 0 & 60 & 200 & $-58$ \\
		خوشه‌ی اولیه & 0 & 120 & 50 & $-46$ \\
		خوشه‌زنی منظم & $-5$ & 100 & 100 & $-46$ \\
		\bottomrule
	\end{tabular}
\end{center}

\begin{keybox}[سازگاری نرخ شلیک، به‌صورت کمّی]
	چون نرخ شلیک در \lr{AdEx} \emph{در طول یک آزمایش} تغییر می‌کند، بهتر است دو منحنی \lr{F-I} تعریف کنیم: نرخ \textbf{لحظه‌ای} (بر اساس اولین \lr{ISI}، پیش از آن‌که $w$ ساخته شود) و نرخ \textbf{ماندگار} (بر اساس \lr{ISI}‌های انتهایی، پس از تعادل یافتن $w$). فاصله‌ی بین این دو منحنی، به‌عنوان تابعی از $I$، خودِ سازگاری نرخ شلیک را آشکار می‌کند. در \lr{LIF} چنین فاصله‌ای وجود ندارد، چون تحت جریان ثابت، \lr{ISI}ها ثابت‌اند.
\end{keybox}

% =====================================================================
\section{تحلیل صفحه‌ی فاز}
% =====================================================================
صفحه‌ی فاز $(v,w)$ به‌صورت هندسی توضیح می‌دهد که \emph{چرا} رژیم‌های بالا پدیدار می‌شوند و آن‌ها را به شانسِ پارامتری فرو نمی‌کاهد.

\begin{defbox}[نال‌کلاین‌ها]
	\textbf{نال‌کلاین $v$}
	($\dot v = 0$)
	و
	\textbf{نال‌کلاین $w$}
	($\dot w = 0$)
	به‌ترتیب چنین‌اند:
	$$w = -g_L(v-E_L) + g_L\Delta_T e^{(v-v_T)/\Delta_T} + I \qquad\text{(نال‌کلاین }v\text{)}$$
	$$w = a(v - E_L) \qquad\text{(نال‌کلاین }w\text{؛ یک خط راست که از }(E_L,0)\text{ می‌گذرد و شیب آن }a\text{ است)}$$
\end{defbox}

\begin{defbox}[نقطه‌ی ثابت]
	هر نقطه‌ی $(v,w)$ که در آن \emph{هر دو} نال‌کلاین همدیگر را قطع کنند، یک نقطه‌ی ثابت است؛ یعنی
	$\dot v = \dot w = 0$.
	اگر مسیر سیستم دقیقاً روی آن نقطه قرار بگیرد، برای همیشه همان‌جا باقی می‌ماند. دور از نقاط ثابت، مسیر سیستم بر اساس دو معادله‌ی مشتق حرکت می‌کند.
\end{defbox}

\subsection{سازوکار آغاز اسپایک}
در $I=0$، سیستم معمولاً \textbf{دو} نقطه‌ی ثابت دارد: یکی پایدار نزدیک $E_L$ (حالت استراحت) و دیگری ناپایدار در موقعیتی جلوتر روی محور $v$ (که نقش یک آستانه‌ی واقعی و هندسی را بازی می‌کند: اگر سیستم از سمت راست آن شروع شود، به‌جای آرام گرفتن، به‌سمت اسپایک رانده می‌شود). با افزایش $I$، نال‌کلاین $v$ به سمت بالا جابه‌جا می‌شود و این دو نقطه‌ی ثابت به هم نزدیک می‌شوند تا جایی که در یک جریان بحرانی با هم \textbf{برخورد کرده و نابود می‌شوند}
(\lr{saddle-node bifurcation}).
وقتی دیگر هیچ نقطه‌ی ثابتی برای استراحت باقی نماند، مسیر سیستم ناچار است به‌طور نامحدود وارد چرخه‌ی شلیک تکراری شود.

\begin{keybox}
	این دقیقاً همان دلیلی است که باعث می‌شود اصلاً رئوبیس برای \lr{EIF} و \lr{AdEx} وجود داشته باشد: رئوبیس جریانی است که در آن نقطه‌ی ثابت پایدار ناپدید می‌شود. بُعد دوم \lr{AdEx} یعنی $w$ همان چیزی است که اجازه می‌دهد همین سازوکار پایه‌ای، سازگاری و \lr{bursting} را نیز پشتیبانی کند: شیب نال‌کلاین $w$ (یعنی $a$) و مکان بازنشانی نسبت به نال‌کلاین‌ها تعیین می‌کنند که مسیر پس از اسپایک کجا فرود بیاید؛ جایی که یا سریعاً یک شبه‌نقطه‌ی ثابت جدید پیدا می‌کند (شلیک تونیک)، یا در هر چرخه دورتر رانده می‌شود (سازگاری)، یا به‌طور دوره‌ای به یک نقطه‌ی ثابت گذرا نزدیک و دوباره از آن دور می‌شود (خوشه‌زنی).
\end{keybox}

% =====================================================================
\section{دینامیک سیناپسی}
% =====================================================================
\subsection{سیناپس نمایی}
\begin{eqbox}[هسته‌ی نمایی]
	$$\tau_{syn}\frac{\dd g}{\dd t} = -g, \qquad g \leftarrow g + w_{syn} \text{ به‌ازای هر اسپایک ورودی}$$
	پاسخ فرم‌بسته به یک ضربه‌ی منفرد در $t=0$:
	$$g(t) = w_{syn}\,e^{-t/\tau_{syn}}$$
\end{eqbox}

\subsection{سیناپس آلفا}
این مدل نسبت به سیناپس نمایی، هسته‌ای با صعود و افت نرم‌تر تولید می‌کند و از دو معادله‌ی دیفرانسیل خطیِ جفت‌شده ساخته می‌شود
(هر دو با یک $\tau$ یکسان):
$$\frac{\dd x}{\dd t} = -\frac{x}{\tau} \quad (x \leftarrow x + w_{syn}\text{ به‌ازای هر اسپایک}), \qquad \frac{\dd g}{\dd t} = \frac{x - g}{\tau}$$
در نتیجه، $g(t)$ دقیقاً در زمان $t=\tau$ پس از یک اسپایک ورودی منفرد به بیشینه می‌رسد
(تابع آلفای \lr{Rall}).

\subsection{ورودی زمینه‌ای پواسونی}
\begin{defbox}[تقریب برنولی در هر بازه‌ی زمانی]
	در هر گام زمانی $\Delta t$ بر حسب میلی‌ثانیه، یک منبع با نرخ $R$ بر حسب هرتز، به‌صورت مستقل با احتمال زیر اسپایک می‌زند:
	$$p_{spike} = \frac{R\cdot\Delta t}{1000}$$
	این تقریب برای یک فرایند پواسون واقعی، هر زمان که
	$R\Delta t \ll 1000$
	باشد، بسیار دقیق است؛ و برای
	$\Delta t \le 0.1$~ms
	و هر نرخ زیستی معقولی کاملاً برقرار است.
\end{defbox}

\begin{keybox}[آمار نویز ضربه‌ای \lr{(shot noise)} — چرا برای تنظیم شبکه مهم است؟]
	اگر یک قطار اسپایک پواسونی با نرخ $R$ (بر حسب اسپایک بر میلی‌ثانیه) را از یک سیناپس نمایی با وزن $w$ و ثابت فروکاهش $\tau$ عبور دهیم، جریان حاصل دارای ویژگی‌های زیر است:
	$$\text{میانگین} = R\,w\,\tau, \qquad \text{واریانس} \approx \frac{R\,w^2\,\tau}{2}$$
	برای یک \emph{میانگین ثابت}، اگر $w$ را زیاد و $R$ را به‌طور متناسب کم کنید، \textbf{واریانس افزایش می‌یابد}؛ چون تعداد رویدادها کمتر اما اندازه‌ی هر رویداد بزرگ‌تر می‌شود. این همان درجه‌ی آزادی‌ای است که شلیک \textbf{میانگین‌محور} را از شلیک \textbf{نوسان‌محور} جدا می‌کند. در حالت میانگین‌محور، میانگین جریان به‌تنهایی آن‌قدر بزرگ است که ولتاژ را به‌طور پیوسته از آستانه رد کند و \lr{ISI}ها نسبتاً منظم باشند. در حالت نوسان‌محور، میانگین زیر رئوبیس نگه داشته می‌شود، اما واریانس آن‌قدر زیاد است که نوسان‌ها گهگاه از آستانه عبور کنند؛ در نتیجه \lr{ISI}ها نامنظم‌تر می‌شوند و $CV$ به 1 نزدیک‌تر است. شبکه‌ی متوازن به همین رژیم نوسان‌محور متکی است.
\end{keybox}

% =====================================================================
\section{شبکه‌های بازگشتی متوازن}
% =====================================================================
اکنون از یک نورون منفرد به $N$ نورون
(در اینجا 1000 نورون)
از نوع \lr{LIF} با اتصال تُنُک می‌رسیم.

\begin{defbox}[قانون دیل \lr{Dale's law}]
	تمام وزن‌های سیناپسیِ \emph{خروجی} یک نورون یا همگی تحریکی‌اند
	(مثبت)
	یا همگی مهاری‌اند
	(منفی)؛
	هرگز ترکیبی از هر دو نیستند. این همان محدودیت زیستیِ واقعی است که یک نورون تنها یک نوع اصلی از ناقل عصبی آزاد می‌کند. در این پروژه، 80\% اول اندیس‌های نورونی تحریکی و 20\% باقی‌مانده مهاری در نظر گرفته شده‌اند.
\end{defbox}

\subsection{ساخت شبکه}
\begin{itemize}[itemsep=2pt]
	\item احتمال اتصال $p$: هر زوج مرتب $(i,j)$ با احتمال مستقل $p$ به هم متصل می‌شود؛ بنابراین میانگین درجه‌ی ورودی تقریباً برابر $N\!\cdot\!p$ است و به‌صورت تقریبی توزیع دوجمله‌ای دارد.
	\item وزن‌های مهاری مقیاس می‌شوند:
	$w_{inh} = -g \cdot w_{exc}$،
	که معمولاً در آن $g$ به‌مراتب بزرگ‌تر از 1 است
	(مثلاً $g=5$)؛
	یعنی مهار باید صرفاً حاضر نباشد، بلکه \emph{در هر اتصال قوی‌تر} هم باشد.
	\item هر نورون علاوه بر این، یک درایو زمینه‌ای پواسونی مستقل نیز دریافت می‌کند
	(بخش \lr{8.3})؛
	بدون این ورودی، یک شبکه‌ی نشت‌دارِ تُنُک به‌سادگی برای همیشه خاموش می‌ماند.
\end{itemize}

\subsection{رژیم نامنظمِ ناهمگام \lr{(Asynchronous-Irregular, AI)}}
\begin{defbox}[رژیم \lr{AI} — سه نشانه‌ی اصلی]
	\begin{enumerate}[itemsep=2pt, leftmargin=1.4em]
		\item \textbf{ناهمگام}: هیچ ریتم سراسریِ مشترک در جمعیت وجود ندارد؛ نمودار رستر ستون‌های عمودی یا خوشه‌های هم‌زمان نشان نمی‌دهد و ردّ نرخ جمعیت تخت است، نه نوسانی.
		\item \textbf{نامنظم}: قطار اسپایکِ نورون‌های منفرد شبیه یک فرایند بی‌حافظه
		(\lr{Poisson-like})
		است. این موضوع با \textbf{ضریب تغییراتِ \lr{ISI}} سنجیده می‌شود:
		$$CV = \frac{\text{std}(\text{ISIs})}{\text{mean}(\text{ISIs})}, \qquad CV\approx 0 \text{ (ساعت‌وار)}, \; CV\approx 1 \text{ (شبیه پواسون)}, \; CV>1 \text{ (خوشه‌ای)}$$
		\item \textbf{پایدار و کم‌نرخ}: نرخ میانگین جمعیت به یک مقدار پایین می‌رسد و همان‌جا می‌ماند، نه این‌که واگرا شود
		(تحریک فراری)
		یا محو شود
		(غلبه‌ی کامل مهار).
	\end{enumerate}
\end{defbox}

\subsection{نرخ جمعیت}
اگر همه‌ی اسپایک‌های شبکه را در پنجره‌هایی با عرض
$\Delta_{bin}$
بر حسب ثانیه دسته‌بندی کنیم، نرخ جمعیت چنین است:
$$r(t) = \frac{\text{تعداد اسپایک در پنجره}}{N \cdot \Delta_{bin}} \quad \text{[Hz]}$$

\subsection{چرا به آن «متوازن» می‌گوییم؟ سازوکار، نه فقط نام}
\begin{keybox}
	وقتی مهار \textbf{ضعیف} باشد
	($g$ کوچک)،
	افزایش وزن تحریکی مستقیماً نرخ جمعیت را بالا می‌برد؛ یعنی بازخورد تحریکیِ بازگشتی تقریباً بدون مهار مؤثر خودش را تقویت می‌کند و شبکه به سمت رفتار فراری یا میانگین‌محور می‌رود. اما وقتی مهار \textbf{قوی} باشد
	($g$ بزرگ)،
	نرخ جمعیت تا حد زیادی نسبت به افزایش بیشترِ وزن تحریکی \emph{حساسیت کمی} نشان می‌دهد، چون مهار تقریباً به همان سرعت که تحریک اضافه می‌رسد، آن را دنبال و خنثی می‌کند. همین کاهش حساسیت — نه یک نرخ عددیِ خاص — تعریف واقعی «تعادل» است. این سازوکار ورودی خالص هر نورون را در رژیم نوسان‌محور نگه می‌دارد
	(جریان خالص کوچک اما پرنویز)
	نه در رژیم میانگین‌محور
	(جریان خالص بزرگ و پایدار)،
	و همین وضعیت است که شلیک نامنظم و ناهمگام توضیح‌داده‌شده در بالا را تولید می‌کند.
\end{keybox}

% =====================================================================
\newpage
\section{Quick-Reference Formula Sheet}
% =====================================================================
\begin{center}
	\renewcommand{\arraystretch}{1.35}
	\begin{longtable}{@{}p{3.6cm}p{12.5cm}@{}}
		\toprule
		\textbf{Concept} & \textbf{Formula} \\
		\midrule
		\endhead
		Passive membrane & $\tau_m \dot v = -(v-v_{rest}) + RI(t)$, \quad $\tau_m = C_m/g_L$ \\
		Passive step response & $v(t) = v_{rest} + \dfrac{I}{g_L}\left(1-e^{-t/\tau_m}\right)$ \\
		LIF rule & $v\!\ge\! v_{th} \Rightarrow v\!\leftarrow\! v_{reset}$, refractory $t_{ref}$ \\
		LIF rheobase & $I_{rheo} = g_L(v_{th}-v_{rest})$ \\
		LIF F-I curve & $f(I) = \left[t_{ref} + \tau_m \ln\!\frac{IR}{IR-(v_{th}-v_{rest})}\right]^{-1}$ \\
		EIF & $C_m \dot v = -g_L(v-E_L) + g_L\Delta_T e^{(v-v_T)/\Delta_T} + I$ \\
		AdEx ($v$) & $C_m \dot v = -g_L(v-E_L) + g_L\Delta_T e^{(v-v_T)/\Delta_T} - w + I$ \\
		AdEx ($w$) & $\tau_w \dot w = a(v-E_L) - w + b\tau_w\sum_f \delta(t-t^f)$ \\
		AdEx spike reset & $v\!\leftarrow\! v_{reset}$, \quad $w \leftarrow w + b$ \\
		v-nullcline & $w = -g_L(v-E_L) + g_L\Delta_T e^{(v-v_T)/\Delta_T} + I$ \\
		w-nullcline & $w = a(v-E_L)$ \\
		Euler step & $x_{n+1} = x_n + \Delta t\, f(x_n)$ \\
		RK4 step & $x_{n+1} = x_n + \frac{\Delta t}{6}(k_1+2k_2+2k_3+k_4)$ \\
		Exponential synapse & $\tau_{syn}\dot g = -g$, \; $g \leftarrow g+w_{syn}$ per spike \\
		Alpha synapse & $\dot x = -x/\tau$ (jump $w_{syn}$ per spike), \; $\dot g = (x-g)/\tau$ \\
		Poisson bin probability & $p_{spike} = R\Delta t / 1000$ \\
		Shot-noise mean/variance & mean $=Rw\tau$, \; variance $\approx Rw^2\tau/2$ \\
		ISI coefficient of variation & $CV = \text{std(ISI)}/\text{mean(ISI)}$ \\
		Population rate & $r(t) = \dfrac{\text{spikes in bin}}{N\cdot\Delta_{bin}\text{[s]}}$ \\
		Dale's law weight scaling & $w_{inh} = -g\cdot w_{exc}$, \; $g \gg 1$ for balance \\
		\bottomrule
	\end{longtable}
\end{center}


\newpage

% =====================================================================
\section{پرسش‌های خودآزمایی}
% =====================================================================
\begin{enumerate}[itemsep=6pt, leftmargin=1.6em]
	\item چرا یک غشای غیرفعال، مستقل از بزرگی جریان ورودی، هرگز نمی‌تواند اسپایک تولید کند؟
	\item توضیح دهید یا به‌دست آورید که چرا نرخ شلیک در منحنی \lr{F-I} مدل \lr{LIF} وقتی $I \to \infty$ می‌رود، به $1/t_{ref}$ نزدیک می‌شود.
	\item چرا \lr{EIF} با داشتن همان $v_T = v_{th}$، رئوبیس \emph{کم‌تری} از \lr{LIF} دارد؟
	\item وقتی $\Delta_T \to 0$، از نظر شهودی چه اتفاقی برای مدل \lr{EIF} می‌افتد؟
	\item جمله‌ی
	$b\,\tau_w\sum_f \delta(t-t^f)$
	در زمان وقوع یک اسپایک دقیقاً چه اثری روی $w$ می‌گذارد؟ فقط در یک جمله و بدون نمادگذاری انتگرالی توضیح دهید.
	\item چرا \lr{AdEx} برای تولید \lr{bursting} به \emph{دو} متغیر حالت نیاز دارد، اما \lr{LIF}/\lr{EIF} با یک متغیر نمی‌توانند این کار را انجام دهند؟
	\item در صفحه‌ی فاز $(v,w)$، از نظر هندسی «در حالت استراحت بودن» برای نورون چه معنایی دارد؟ و وقتی جریان ورودی سیستم را به سمت رئوبیس می‌برد، چه رخ می‌دهد؟
	\item تفاوت بین منحنی \lr{F-I} «لحظه‌ای» و «ماندگار» چیست، و چرا فقط \lr{AdEx} — نه \lr{LIF} — بین این دو فاصله دارد؟
	\item چرا یک منبع زمینه‌ای پواسونی، برای آن‌که واقعاً شلیک ایجاد کند، هم به نرخ ناصفر \emph{و} هم به وزن سیناپسی کافی نیاز دارد، نه فقط به نرخ بالا؟
	\item تفاوت رژیم شلیک «میانگین‌محور» و «نوسان‌محور» را با زبان میانگین و واریانس نویز ضربه‌ای توضیح دهید.
	\item قانون دیل دقیقاً چه محدودیتی روی ماتریس اتصال $W$ اعمال می‌کند؟
	\item چرا رژیم نامنظمِ ناهمگام \lr{(AI)} یک \emph{ویژگی برآمده از شبکه} محسوب می‌شود، نه ویژگی یک نورون منفرد؟
	\item اگر در یک شبکه‌ی متوازن، همه‌ی وزن‌های تحریکی را دو برابر کنید \emph{بدون} آن‌که مهار را به همان نسبت افزایش دهید، انتظار دارید چه بر سر نرخ جمعیت بیاید؟ و اگر پیش از این افزایش، مهار از قبل بسیار قوی بوده باشد چه؟
	\item چرا این پروژه برای تمام شبیه‌سازی‌های شبکه‌ای از اویلر استفاده می‌کند، اما \lr{RK4} را فقط برای نورون‌های منفرد نگه می‌دارد؟
\end{enumerate}

\vspace{0.3cm}
\begin{center}\rule{0.4\textwidth}{0.4pt}\end{center}

\newpage

\subsection*{پاسخ‌های کوتاه}
{\small
	\begin{enumerate}[itemsep=4pt, leftmargin=1.6em]
		\item چون در مدل اصلاً هیچ قانون آستانه/بازنشانی وجود ندارد؛ $v$ فقط به‌صورت مجانبی به سمت $v_{rest}+IR$ می‌رود و هیچ سازوکاری برای فعال‌کردن یک رویداد اسپایک وجود ندارد.
		\item چون وقتی $I\to\infty$، زمان شارژ از بازنشانی تا آستانه به صفر میل می‌کند
		(جمله‌ی لگاریتمی از بین می‌رود)
		و تنها فاصله‌ی باقی‌مانده بین اسپایک‌ها همان دوره‌ی نارسانایی است؛ بنابراین
		$f \to 1/t_{ref}$.
		\item چون جریان نمایی حتی زیر $v_T$ هم ناصفر و دپولاریزان است و بخشی از کاری را که در آستانه‌ی سخت نیازمند جریان بیشتر بود، از قبل انجام می‌دهد.
		\item جمله‌ی نمایی را تیزتر می‌کند تا به چیزی شبیه یک پله در $v_T$ تبدیل شود و \lr{EIF} هرچه بیشتر به \lr{LIF} شبیه شود.
		\item در هر اسپایک، $w$ فوراً دقیقاً به‌اندازه‌ی $b$ بالا می‌پرد؛ بین اسپایک‌ها این جمله هیچ سهمی ندارد و $w$ فقط به سمت $a(v-E_L)$ آرام می‌گیرد.
		\item یک متغیر منفرد بعد از هر اسپایک به همان وضعیت بازنشانی می‌شود، بنابراین در ورودی ثابت همه‌ی \lr{ISI}ها یکسان می‌شوند؛ اما $w$ حافظه‌ی شلیک‌های اخیر را بین فاصله‌های زمانی حمل می‌کند و اجازه می‌دهد \lr{ISI}های پیاپی تغییر کنند، بلندتر یا کوتاه‌تر شوند، یا خوشه بسازند.
		\item استراحت یعنی قرار گرفتن روی یک نقطه‌ی ثابت پایدار
		(تقاطع دو نال‌کلاین و جهت میدان به‌سمت داخل).
		با افزایش $I$، نال‌کلاین $v$ جابه‌جا می‌شود و زوجِ پایدار/ناپایدارِ نقطه‌ی ثابت به هم می‌رسند و ناپدید می‌شوند
		(\lr{saddle-node bifurcation})؛
		از آن لحظه به بعد دیگر جای استراحتی وجود ندارد و سیستم ناچار وارد چرخه‌ی شلیک تکراری می‌شود.
		\item لحظه‌ای یعنی نرخ به‌دست‌آمده از اولین \lr{ISI}، پیش از آن‌که $w$ ساخته شود؛ ماندگار یعنی نرخ به‌دست‌آمده از \lr{ISI}های انتهایی، بعد از تعادل $w$. در \lr{LIF} متغیر سازگاری وجود ندارد، پس نرخ در طول آزمایش تغییر نمی‌کند و فاصله‌ای بین این دو منحنی دیده نمی‌شود.
		\item نرخ فقط تعیین می‌کند ضربه‌ها \emph{چند وقت یک‌بار} برسند؛ وزن تعیین می‌کند هر ضربه \emph{چقدر بزرگ} باشد. نرخ بالای ضربه‌های بسیار کوچک ممکن است هنوز هم به دپولاریزاسیون خالص مؤثری نینجامد، اگر وزن آن‌قدر کم باشد که چه به‌صورت منفرد و چه در مجموع نسبت به آستانه بی‌اثر بماند.
		\item در حالت میانگین‌محور، میانگین جریان به‌تنهایی تقریباً در هر چرخه $v$ را از آستانه رد می‌کند و \lr{ISI}ها نسبتاً منظم‌اند. در حالت نوسان‌محور، میانگین زیر آستانه می‌ماند و فقط نوسان‌های تصادفی
		(واریانس)
		گهگاه نورون را از آستانه عبور می‌دهند؛ در نتیجه \lr{ISI}ها نامنظم و شبیه پواسون می‌شوند.
		\item هر \emph{ستون} از $W$
		(یعنی تمام وزن‌های خروجی یک نورون پیش‌سیناپسی)
		باید یا کاملاً
		$\ge 0$
		باشد
		(منبع تحریکی)
		یا کاملاً
		$\le 0$
		باشد
		(منبع مهاری)؛
		در یک ستون نمی‌توان علامت‌های آمیخته داشت.
		\item هیچ نورون \lr{LIF} منفردی در شبکه، از نظر کیفی رفتاری متفاوت از یک نورون \lr{LIF} منزوی با ورودی نویزی ندارد؛ رفتار جمعیتیِ ناهمگام، نامنظم و پایدار فقط از نحوه‌ی سیم‌کشی 1000 نورون با هم پدید می‌آید
		(قانون دیل + اتصال تُنُک + مهارِ قوی‌تر از تحریک).
		\item اگر مهار در ابتدا ضعیف باشد، نرخ به‌طور محسوس بالا می‌رود
		(تقویت بازگشتی).
		اگر مهار از قبل بسیار قوی باشد، نرخ تقریباً ثابت می‌ماند، چون مهار درایو اضافی را جذب می‌کند؛ و این دقیقاً همان سازوکار تعادل است.
		\item چون شبکه‌ها شامل نورون‌های زیاد و گام‌های زمانی فراوان‌اند و هزینه‌ی اویلر
		(یک ارزیابی مشتق در هر گام)
		است که شبیه‌سازی را از نظر محاسباتی ممکن می‌کند؛ هزینه‌ی $4\times$ بیشتر \lr{RK4} فقط زمانی می‌ارزد که بخواهیم مسیر یک نورون منفرد را با دقت بالا در برابر یک مرجع اعتبارسنجی کنیم.
	\end{enumerate}
}

\newpage

% =====================================================================
\section{واژه‌نامه}
% =====================================================================
\begin{description}[style=nextline, leftmargin=1.2em, itemsep=3pt]
	\item[رئوبیس] کمینه‌ی جریان ثابت لازم برای شلیکِ پایدار یا نهایتاً رخ‌دادنی.
	\item[دوره‌ی نارسانایی ($t_{ref}$)] بازه‌ی زمانی پس از یک اسپایک که در آن نورون نمی‌تواند دوباره شلیک کند؛ بیشینه‌ی نرخ شلیک ممکن را به $1/t_{ref}$ محدود می‌کند.
	\item[نال‌کلاین] منحنی‌ای در صفحه‌ی فاز که روی آن مشتق یکی از متغیرهای حالت صفر است.
	\item[نقطه‌ی ثابت] نقطه‌ای که در آن مشتق \emph{همه‌ی} متغیرهای حالت به‌طور هم‌زمان صفر است.
	\item[دوشاخه‌گی \lr{(Bifurcation)}] تغییر کیفی در رفتار یک سیستم دینامیکی
	(مثلاً تعداد نقاط ثابت)
	وقتی یک پارامتر از یک مقدار بحرانی عبور می‌کند.
	\item[دوشاخه‌گی گره--زین \lr{(Saddle-node bifurcation)}] وضعیتی که در آن دو نقطه‌ی ثابت
	(یکی پایدار و دیگری ناپایدار)
	با تغییر یک پارامتر به هم برخورد کرده و نابود می‌شوند؛ سازوکار آغاز اسپایک در \lr{EIF}/\lr{AdEx}.
	\item[تحریک‌پذیری نوع اوّل] حالتی که در آن نرخ شلیک در رئوبیس به‌صورت پیوسته از صفر افزایش می‌یابد
	(بر خلاف نوع دوم که به یک نرخ ناصفر می‌جهد).
	\item[سازگاری نرخ شلیک \lr{(SFA)}] بلندتر شدن تدریجی \lr{ISI}ها در طول یک قطار اسپایک تحت ورودی ثابت.
	\item[قانون دیل] سیناپس‌های خروجی یک نورون یا همگی تحریکی‌اند یا همگی مهاری؛ ترکیبی از هر دو مجاز نیست.
	\item[ضریب تغییرات \lr{(CV)}] نسبت انحراف معیار به میانگینِ فاصله‌های بین اسپایکی؛ یک معیار استاندارد برای نامنظمی.
	\item[رژیم نامنظمِ ناهمگام \lr{(AI)}] حالت شبکه‌ای بدون ریتم مشترک، با شلیک‌های منفرد نامنظم و نرخ میانگین پایدار و کم.
	\item[میانگین‌محور در برابر نوسان‌محور] این‌که آیا شلیک عمدتاً به‌واسطه‌ی عبورِ خودِ جریان میانگین از آستانه رخ می‌دهد، یا به‌واسطه‌ی نویز/واریانس که گهگاه یک جریانِ در میانگین زیرآستانه‌ای را از آستانه عبور می‌دهد.
	\item[نویز ضربه‌ای \lr{(Shot noise)}] نوسان‌های جریانی حاصل از عبور دادن یک قطار اسپایک گسسته
	(مثلاً پواسونی)
	از یک هسته‌ی سیناپسی پیوسته.
\end{description}

\end{document}
